\begin{tfjmenum}

Pour cette réponse,
Soit une droite d d'équation $y = kx$ passant par $O(0,0)\in d$,
Un point $F(x',y')\in A$ tel que $x'$ et $y'$ sont premiers entre eux, i.e. F est le premier point de d de coordonnées entières naturelles non nulles. Donc F est le premier arbre sur d.
Et $I(a,b)\in d$ tel que $(a,b)\in\mathbb{R}^2 \setminus A$.

\item

CONDITION NECESSAIRE

Si $F(x',y')\in d$ et $I(a,b)\in d$ tel que $F\in [OI]$ alors Ines est cachée, d'après l'énoncé.

$F(x',y')\in d$ et $I(a,b)\in d$ tel que $F\in [OI]$

$\Leftrightarrow$ $k\in \mathbb{Q}^+$ car $F(x',y')\in A$ et $\vec{OF}$ et $\vec{OI}$ sont colinéaires car $F\in [OI]\subset d$

$\Leftrightarrow$ $k\in \mathbb{Q}^+$ et $\vec{OI}=n\vec{OF}$ avec $n>1$

$\Leftrightarrow$ $k\in \mathbb{Q}^+$ et $\vec{OI}=\begin{pmatrix} a \\ b \end{pmatrix} = n\vec{OF}=\begin{pmatrix} x' \\ y' \end{pmatrix}$ où $n>1$ car $O(0,0)$.

$\Leftrightarrow$ $k\in \mathbb{Q}^+$ et $n=\frac{a}{x'}=\frac{b}{y'}>1$

On a donc :
Si $k\in \mathbb{Q}^+$ et $n=\frac{a}{x'}=\frac{b}{y'}>1$ $(n\neq 0)$ alors Ines est cachée.

Ainsi par contraposée on peut donc conclure :

Si Ines est visible alors $k\in \mathbb{R}\setminus \mathbb{Q}^+$ ou $k\in \mathbb{Q}^+$ et $n=\frac{a}{x'}=\frac{b}{y'}<1$ $(n\neq 0)$

\item

CONDITION SUFFISANTE

\begin{itemize}
    \item Si $k\in \mathbb{R} \setminus \mathbb{Q}^+$ alors il n'existe pas de point $E\in A$ tel que $E\in d$.

    Donc il n'existe pas de point $F(x',y')\in d$ tel que $F(x',y')\in A$.

    Or $I(a,b)\in d$.

    Donc il n'existe pas d'arbre sur le segment $[OI]$.

    Donc Ines est visible.

    \item $k\in \mathbb{Q}^+$ et $n=\frac{a}{x'}=\frac{b}{y'}<1$ $(n\neq 0)$ $\Leftrightarrow$ $k\in \mathbb{Q}^+$ et $\vec{OI}=\begin{pmatrix} a \\ b \end{pmatrix} = n\vec{OF}=\begin{pmatrix} x' \\ y' \end{pmatrix}$ $(n\neq 0)$.

    \begin{itemize}
        \item Si $n<0$ :

        Donc $n=\frac{a}{x'}=\frac{b}{y'}<0$ or $(x',y')\in A$ donc $a<0$ et $b<0$.

        Donc $(a,b)\in \mathbb{R}^{-2}$ or il n'y a pas d'arbre car $A=\mathbb{N}^2 \setminus \{(0,0)\}$.

        Donc Ines est visible.

        \item Si $0<n<1$ :

        Donc $0<n=\frac{a}{x'}=\frac{b}{y'}<1$ donc $a<x'$ et $b<y'$.

        Or $I(a,b)\in d$ et $F(x',y')\in d$.

        Donc $F\notin[OI]$.

        Donc il n'y a pas d'arbre sur $[OI]$, F étant le premier arbre sur d.

        Donc Ines est visible.
    \end{itemize}
\end{itemize}

Conclusion :

Si $k\in \mathbb{R} \setminus \mathbb{Q}^+$ ou $k\in \mathbb{Q}^+$ et $n=\frac{a}{x'}=\frac{b}{y'}<1$ $(n\neq 0)$

alors Ines est visible.

\item

CONDITION NÉCESSAIRE ET SUFFISANTE

On a montré l'implication et la réciproque donc on a :

Ines est visible si et seulement si $k\in \mathbb{Q}^+$ et $n=\frac{a}{x'}=\frac{b}{y'}<1$ $(n\neq 0)$ ou $k\in \mathbb{R} \setminus \mathbb{Q}^+$

\end{tfjmenum}
